\documentclass{getwriting}
\usepackage{amsmath}

\title{A practical guide to reproducible field studies}
\author[1]{Maya Rivera\thanks{maya.rivera@example.edu}\,\orcidlink{0000-0002-1825-0097}}
\author[2]{Jon Bell\thanks{jon.bell@example.org}}
\affil[1]{Centre for Environmental Data, Northbridge University}
\affil[2]{Open Methods Lab, Example Institute}

\begin{document}
\maketitle

\begin{abstract}
This starter shows a readable, double-spaced article workflow for a
format-free submission.  It includes an abstract, keywords, a methods
checklist, an editable equation, a native table, a code listing, and a
plain BibTeX bibliography without depending on external figures.
\end{abstract}

\noindent\textbf{Keywords:} reproducibility, field studies, data quality,
open methods, article template

\section{Introduction}
Field studies often produce useful measurements before a team has settled on
a permanent data model.  A small, explicit record of the sampling plan makes
those measurements easier to audit and reuse.  This article uses a fictional
coastal survey to demonstrate the structure of a manuscript without carrying
over a real author's data or prose.

\section{Study design}
The survey visits three sites at two times of day.  Each observation records
the site identifier, the instrument configuration, and a quality flag.  The
design follows the simple principle that a result should retain enough context
for another researcher to reproduce its collection~\cite{wilkinson2016}.

\begin{table}[H]
  \centering
  \caption{A compact sampling register.}
  \label{tab:register}
  \begin{tabular}{llll}
    \toprule
    Site & Morning & Afternoon & Quality rule \\
    \midrule
    North inlet & 12 & 12 & sensor check \\
    Estuary bend & 10 & 10 & duplicate 10\% \\
    South marsh & 8 & 8 & field blank \\
    \bottomrule
  \end{tabular}
\end{table}

\section{Analysis plan}
Let $x_i$ be an observation and $q_i$ its quality weight.  The weighted mean
reported for a site is
\begin{equation}
  \bar{x}_q = \frac{\sum_i q_i x_i}{\sum_i q_i},
  \label{eq:weighted-mean}
\end{equation}
which keeps the calculation visible and makes the unit of analysis explicit.
The uncertainty interval can then be reported alongside the point estimate
rather than hidden in a downstream spreadsheet.

The analysis record is intentionally small:
\begin{enumerate}
  \item freeze the raw observation file and record its checksum;
  \item apply the quality rule in Table~\ref{tab:register};
  \item calculate Eq.~\eqref{eq:weighted-mean};
  \item publish the code and a machine-readable result table.
\end{enumerate}

\section{Reproducibility notes}
The following command is a placeholder for the project-specific analysis
entry point; replace it with the command that recreates the figures and
tables from the frozen data.
\begin{verbatim}
make report DATA=data/observations.csv OUTPUT=build/results.csv
\end{verbatim}
Keeping a short command block in the manuscript helps a reader distinguish
the analytical recipe from the surrounding narrative.

\section{Conclusion}
The \texttt{getwriting} class is deliberately close to the standard article
class.  Start by replacing the fictional study, then keep or remove the
table, equation, and code block as the paper requires.

\bibliography{references}
\end{document}
